\section*{\texorpdfstring{\S\CurrentExerciseGroup}{Section \CurrentExerciseGroup}}
\phantomsection
\addcontentsline{toc}{section}{\texorpdfstring{Exercises for \S\CurrentExerciseGroup}{Exercises for Section \CurrentExerciseGroup}}

\begin{enumerate}
\def\labelenumi{\arabic{enumi})}
\tightlist
\item
  Let E be a Riesz space, x an element \textgreater0 of E. If there exists a relatively bounded linear form L on E such that \(L(x) \neq 0\) , show that the set of nx (n an integer \textgreater0) cannot be bounded above in E. In particular, show that there does not exist any relatively bounded linear form on the Riesz space E defined in Exer. 5 b) of §1.
\end{enumerate}

¶ 2) a) Let L be a positive linear form on the Riesz space E. Consider an element a \textgreater{} 0 of E and the set of positive linear forms \(M \leqslant L\) such that: \(1^{\circ} M(x) = L(x)\) for every x such that \(0 \leqslant x \leqslant a\) ; \(2^{\circ} M(x) = 0\) for every \(x \geqslant 0\) alien to a. Show that this set of positive linear forms has a largest element \(L_{a}\) and that, for every \(x \geqslant 0\) , \(L_{a}(x) = \inf L(y)\) , where y runs over the set of all elements such that \(0 \leqslant y \leqslant x\) and x - y is alien to a. Show that \(L_{\lambda a} = L_{a}\) for every scalar \(\lambda > 0\) , and that if a and b are alien then \(L_{a+b} \leqslant L_{a} + L_{b}\) . If E is fully lattice-ordered, show that if a and b are alien in E then \(L_{a+b} = L_{a} + L_{b}\) .

\begin{enumerate}
\def\labelenumi{\alph{enumi})}
\setcounter{enumi}{1}
\tightlist
\item
  Take E to be the Riesz space of continuous real-valued functions on the interval \(\mathbf{I} = [0,1]\) of \(\mathbf{R}\) , and \(L(x) = x\left(\frac{1}{2}\right)\) ; show by means of an example that one can have \(L_{a + b} < L_a + L_b\) for two alien elements \(a, b\) of \(\mathbf{E}\) .
\end{enumerate}

\begin{enumerate}
\def\labelenumi{\arabic{enumi})}
\setcounter{enumi}{2}
\tightlist
\item
  Let \(\mathbf{E}\) be a fully lattice-ordered space.
\end{enumerate}

\begin{enumerate}
\def\labelenumi{\alph{enumi})}
\item
  Show that every linear form \(L\) on \(\mathbf{E}\) that is continuous for the topology \(\mathcal{T}_0(\mathbf{E})\) (§1, Exer. 10) is relatively bounded, and that \(|L|\) is continuous for \(\mathcal{T}_0(\mathbf{E})\) (argue by contradiction, on observing that for every \(x \in \mathbf{E}\) the sequence of elements \(x / n\) tends to 0 as \(n\) tends to infinity, for the topology \(\mathcal{T}_0(\mathbf{E})\) ).
\item
  Let A be any infinite set, U an ultrafilter on A whose sets have empty intersection. On the fully lattice-ordered space \(E = \mathcal{B}(A)\) , consider the positive linear form \(x \mapsto \lim_{\mu} x(t)\) ; show that this linear form is not continuous for the topology \(\mathcal{T}_{0}(E)\) .
\end{enumerate}

\begin{enumerate}
\def\labelenumi{\arabic{enumi})}
\setcounter{enumi}{3}
\tightlist
\item
  Let E be a fully lattice-ordered space.
\end{enumerate}

\begin{enumerate}
\def\labelenumi{\alph{enumi})}
\item
  Let \(L\) be a positive linear form on \(\mathbf{E}\) , continuous for the topology \(\mathcal{T}_0(\mathbf{E})\) (§1, Exer. 10). Show that the set of \(x \in \mathbf{E}\) such that \(L(|x|) = 0\) is a band \(Z(L)\) . Let \(S(L)\) be the band supplementary to \(Z(L)\) in \(\mathbf{E}\) (§1, Exer. 6).
\item
  Let L and M be two positive linear forms on E, continuous for \(\mathcal{T}_{0}(\mathrm{E})\) . In order that, in the space \(\Omega\) of relatively bounded linear forms on E, M belong to the band generated by L, it is necessary and sufficient that \(S(M) \subset S(L)\) . (To see that the condition is sufficient, argue by contradiction, on supposing that M does not verify the condition of Prop. 5; consider a sequence \((y_{n})\) of elements of E such that \(0 \leqslant y_{n} \leqslant x\) ,
\end{enumerate}

\(L(y_{n}) \leqslant 1 / 2^{n}\) and \(M(y_{n}) \geqslant \alpha > 0\) , and infer the existence of an element \(z \geqslant 0\) in E such that \(L(z) = 0\) and \(M(z) \geqslant \alpha\) .

\begin{enumerate}
\def\labelenumi{\alph{enumi})}
\setcounter{enumi}{2}
\tightlist
\item
  Let \(L\) and \(M\) be two positive linear forms on \(\mathbf{E}\) , continuous for \(\mathcal{T}_0(\mathbf{E})\) . In order that, in \(\Omega\) , \(L\) and \(M\) be alien, it is necessary and sufficient that \(S(L) \cap S(M) = \{0\}\) . (To see that the condition is necessary, prove that if \(S(L) \cap S(M)\) does not reduce to 0 then the bands in \(\Omega\) generated by \(L\) and \(M\) have an element \(\neq 0\) in common, by using b).)
\end{enumerate}

¶ 5) a) Let E be a Riesz space, H an isolated linear subspace of E (§1, Exer. 4). Let Ω be the space of relatively bounded linear forms on E and let Θ be the subspace of Ω consisting of the linear forms \(L \in \Omega\) that are zero on H. Show that Θ is an isolated subspace of Ω and that Θ is a fully lattice-ordered space isomorphic to the space of relatively bounded linear forms on the Riesz space E/H.

\begin{enumerate}
\def\labelenumi{\alph{enumi})}
\setcounter{enumi}{1}
\tightlist
\item
  For a linear mapping \(u\) of \(E\) into a Riesz space \(F\) , the following conditions are equivalent:
\end{enumerate}

\(\alpha) u(\sup(x, y)) = \sup(u(x), u(y));\)

\(\beta) u(\inf(x, y)) = \inf(u(x), u(y));\)

\(\gamma) u(x^{+}) = (u(x))^{+}\) ;

\(\delta)\) the relation \(\inf (x,y) = 0\) implies \(\inf (u(x),u(y)) = 0\)

One then says that u is a latticial linear mapping.

\begin{enumerate}
\def\labelenumi{\alph{enumi})}
\setcounter{enumi}{2}
\item
  For a linear subspace H of E to be maximal in the set of isolated subspaces \(\neq\) E, it is necessary and sufficient that it be of the form \(\operatorname{Ker}(f)\) , where \(f\) is a latticial linear form \(\neq 0\) . (Make use of GT, V, §3, Exer. 1.)
\item
  Give an example of a Riesz space not reduced to 0 that contains no maximal isolated subspace (cf.~§1, Exer. 5 b)).
\end{enumerate}

¶ 6) Let E be a Riesz space, Ω the fully lattice-ordered space of relatively bounded linear forms on E, and F a co-lattice subspace of Ω (§1, Exer. 3). For every \(x \in E\) , the mapping \(x' \mapsto \langle x, x' \rangle\) of F into R is a relatively bounded linear form \(u_x\) on F, and the mapping \(x \mapsto u_x\) is an increasing linear mapping of E into the fully lattice-ordered space Ω' of relatively bounded linear forms on F. In order that \(x \mapsto u_x\) be an isomorphism of E onto a co-lattice subspace of Ω', that is, for \(u_x > 0\) to imply x \textgreater{} 0, and that \(u_{\text{sup}(x,y)} = \text{sup}(u_x, u_y)\) , it is necessary and sufficient that the following two conditions be verified: 1° for every x \textgreater{} 0 in E, there exists an \(x' > 0\) in F such that \(\langle x, x' \rangle > 0\) ; 2° for every pair of alien elements \(y \geqslant 0\) , \(z \geqslant 0\) of E, for every number \(\varepsilon > 0\) and for every \(x' \geqslant 0\) in F, there exist two elements \(y' \geqslant 0\) , \(z' \geqslant 0\) of F such that \(x' = y' + z'\) and \(\langle y, y' \rangle + \langle z, z' \rangle \leqslant \varepsilon\) . (Observe that if the second condition is verified and if v is a positive linear form on F such that \(v \geqslant u_x\) , then \(v(x') \geqslant \langle x^+, x' \rangle - \varepsilon\) for every \(x' \geqslant 0\) in F and every \(\varepsilon > 0\) .)

If the condition \(2^{\circ}\) is fulfilled but not condition \(1^{\circ}\) , then the set of \(x \in E\) such that \(\langle x, x' \rangle = 0\) for all \(x' \in F\) is an isolated linear subspace \(H\) of \(E\) . By passage to the quotient, the mapping \(x \mapsto u_x\) then defines an isomorphism of the Riesz space \(E / H\) onto a co-lattice subspace of \(\Omega'\) .

Show that when \(\mathbf{F} = \Omega\) , the above condition \(2^{\circ}\) is always verified (make use of Exer. 2).

\begin{enumerate}
\def\labelenumi{\arabic{enumi})}
\setcounter{enumi}{6}
\tightlist
\item
  Let \(E\) be a Riesz space of finite dimension \(n\) .
\end{enumerate}

\begin{enumerate}
\def\labelenumi{\alph{enumi})}
\item
  If E is archimedean (§1, Exer. 5), show that the cone P of elements \(\geqslant 0\) of E is closed and has an interior point. From this, deduce that the intersection of the support hyperplanes of P reduces to 0 (make use of the fact that \(P \cap (-P) = \{0\}\) ).
\item
  Using a) and Exer. 6, show that every archimedean Riesz space of dimension n is isomorphic to the product space \(R^{n}\) (cf.~§1, Exer. 13 e)).
\item
  Give an example of a totally ordered Riesz space of dimension 2.
\end{enumerate}

¶ 8) Let E be a Riesz space, \((U_{\iota})\) a family of positive linear forms on E. We consider the topology \(\mathcal{T}\) on E defined by the semi-norms \(U_{\iota}(|x|)\) .

\begin{enumerate}
\def\labelenumi{\alph{enumi})}
\item
  Show that the subspace H of E formed by the x such that \(U_{\iota}(|x|)=0\) for all \(\iota\) is an isolated subspace of E (§1, Exer. 4); the Hausdorff space associated with E is the quotient space E/H; by passage to the quotient, the \(U_{\iota}\) define positive linear forms \(\dot{U}_{\iota}\) on E/H, and the quotient topology on E/H is defined by the semi-norms \(\dot{U}_{\iota}(|\dot{x}|)\) .
\item
  Show that in the space E the mapping \(x \mapsto |x|\) is uniformly continuous, and deduce therefrom that if the topology T is Hausdorff then it is compatible with the ordered vector space structure of E.
\item
  Assume that the topology \(\mathcal{T}\) is Hausdorff and that E is fully lattice-ordered; show that every band B in E is closed, and that if B' is the band formed by the elements alien to every element of B, then E is the topological direct sum of B and B'.
\item
  Assume that the topology \(\mathcal{T}\) is Hausdorff. Let \(\overline{\mathrm{E}}\) be the completion of \(\mathrm{E}\) ; if \(\mathrm{P}\) is the set of elements \(\geqslant 0\) of \(\mathrm{E}\) , show that the closure \(\overline{\mathrm{P}}\) of \(\mathrm{P}\) in \(\widehat{\mathrm{E}}\) defines a Riesz space structure on \(\widehat{\mathrm{E}}\) (cf.~§1, No.~2).
\item
  If E is Hausdorff and complete for the topology T, show that in order for a set \(A \subset E\) , directed for the relation \(\leqslant\) , to have a supremum in E, it is necessary and sufficient that, for every index \(\iota\) , \(U_{\iota}(x)\) be bounded above in A; for every continuous, increasing numerical function f on E, one then has \(\sup_{x \in A} f(x) = f(\sup A)\) . Deduce from
\end{enumerate}

this that E is fully lattice-ordered.

\begin{enumerate}
\def\labelenumi{\alph{enumi})}
\setcounter{enumi}{5}
\tightlist
\item
  Assume that E is Hausdorff and complete for the topology T. Show that if a filter F on E has a limit for the order structure of E (§1, Exer. 9), then it converges to the same limit for the topology T (make use of e)).
\end{enumerate}

\begin{enumerate}
\def\labelenumi{\arabic{enumi})}
\setcounter{enumi}{8}
\tightlist
\item
  Let E be a Riesz space, \(\Omega\) the space of relatively bounded linear forms on E. Consider the topology on \(\Omega\) defined by the semi-norms \(L \mapsto |L|(x)\) , where x runs over the set of elements \(\geqslant 0\) of E. Show that \(\Omega\) , equipped with this topology, is Hausdorff and complete.
\end{enumerate}

